%% ─────────────────────────────────────────────────────────────────────────────
%% CV — Talvin Ackbaraly (English Version)
%% Serif (XCharter), black text, thin rules, bullets, no colours/icons
%% Compile: latexmk -pdf cv-talvin-ackbaraly-en.tex
%% ─────────────────────────────────────────────────────────────────────────────
\documentclass[10.5pt, a4paper]{extarticle}

\usepackage[T1]{fontenc}
\usepackage[utf8]{inputenc}
\usepackage[english]{babel}
\usepackage[top=1.3cm, bottom=1.3cm, left=1.8cm, right=1.8cm]{geometry}
\usepackage{XCharter}
\usepackage{microtype}
\usepackage{titlesec}
\usepackage{enumitem}
\usepackage{tabularx}
\usepackage{xcolor}
\usepackage[hidelinks]{hyperref}

\hypersetup{pdftitle={Resume — Talvin Ackbaraly}, pdfauthor={Talvin Ackbaraly}}

\pagestyle{empty}
\setlength{\parindent}{0pt}
\setlength{\parskip}{0pt}

%% Section: small caps, thin rule underneath
\titleformat{\section}{\large\scshape}{}{0pt}{}[\vspace{-6pt}\rule{\textwidth}{0.4pt}]
\titlespacing*{\section}{0pt}{10pt}{5pt}

\setlist[itemize]{leftmargin=1.1em, label={\textbullet}, topsep=2pt, itemsep=1pt, parsep=0pt}

%% \entry{Title}{Date}{Subtitle}{Right-hand detail}
\newcommand{\entry}[4]{%
  \textbf{#1} \hfill #2\\
  \textit{#3} \hfill \textit{#4}\par
}
\newcommand{\tech}[1]{\hfill{\small\textit{#1}}}

\begin{document}

%% ══ HEADER ═══════════════════════════════════════════════════════════════════
\begin{center}
  {\fontsize{24}{28}\selectfont\scshape Talvin Ackbaraly}\\[5pt]
  Engineering student — Computer Vision \& GPU Computing\\[3pt]
  \small
  \href{mailto:talvin.ackbaraly@icloud.com}{talvin.ackbaraly@icloud.com} \enspace·\enspace
  \href{tel:+33769389008}{+33 7 69 38 90 08} \enspace·\enspace
  \href{https://talvin-ackbaraly.com}{talvin-ackbaraly.com}\\
  \href{https://github.com/talvinckb}{github.com/talvinckb} \enspace·\enspace
  \href{https://www.linkedin.com/in/talvin-ackbaraly}{linkedin.com/in/talvin-ackbaraly}\\[3pt]
  \textit{Seeking a 6-month final-year internship starting February 2027}
\end{center}

%% ══ EDUCATION ════════════════════════════════════════════════════════════════
\section{Education}

\entry{EPITA — Engineering School of Computer Science}{Sept. 2022 -- 2027}
      {Master of Engineering, Image specialization}{Paris, France}
{\small Image processing \& synthesis, computer vision, deep learning, GPU computing.}

\vspace{4pt}
\entry{Université du Québec à Chicoutimi (UQAC)}{Jan. -- May 2024}
      {Academic exchange semester}{Québec, Canada}

%% ══ PROJECTS ═════════════════════════════════════════════════════════════════
\section{Projects}

\textbf{Illustration Detection for the BnF} — \textit{final-year project, team of 4, ongoing} \tech{PyTorch, YOLO26, Florence-2}
\begin{itemize}
  \item Building a pipeline to \textbf{detect and classify illustrations} in digitized documents of the French National Library.
  \item Training and benchmarking detection and classification models (YOLO26, Florence-2) on a \textbf{30\,GB dataset}.
  \item Defined the evaluation metrics with the client and present results in weekly meetings.
\end{itemize}

\vspace{3pt}
\textbf{Real-Time Motion Detection on GPU} — \textit{team of 4} \tech{CUDA, C++, Nsight}
\begin{itemize}
  \item Ported a video motion-detection pipeline from CPU to GPU, achieving a \textbf{24$\times$ speedup}.
  \item Profiled and removed bottlenecks with NVIDIA Nsight.
\end{itemize}

\vspace{3pt}
\textbf{Real-Time Fluid Simulation} — \textit{team of 2} \tech{C++, OpenGL, GLSL}
\begin{itemize}
  \item Built an SPH fluid engine simulating \textbf{75,000 particles at 60+ FPS}, with physics in compute shaders.
\end{itemize}

\vspace{3pt}
\textbf{Automatic License Plate Recognition} — \textit{team of 2} \tech{C++, Python, OpenCV}
\begin{itemize}
  \item Built the full pipeline (preprocessing, candidate extraction, HOG features, Random Forest): \textbf{86\% precision}.
  \item Benchmarked Python and C++ implementations; ran the project through GitLab issues, MRs and CI.
\end{itemize}

\vspace{3pt}
\textbf{Lung Function Prediction from CT Scans} — \textit{team of 4} \tech{Python, XGBoost, React, Docker}
\begin{itemize}
  \item Predicted forced vital capacity from lung segmentation using classical ML only (XGBoost, no CNNs allowed).
  \item Delivered it as a web app with a React front-end, containerized with Docker and a GitLab CI pipeline.
\end{itemize}

\vspace{3pt}
\textbf{Systems Programming} — \textit{EPITA projects} \tech{C, Linux}
\begin{itemize}
  \item Implemented a memory allocator (malloc), a POSIX shell (42sh), a BitTorrent client and a \texttt{find} clone.
\end{itemize}

%% ══ EXPERIENCE ═══════════════════════════════════════════════════════════════
\section{Experience}

\entry{Studiocanal}{Sept. 2025 -- Jan. 2026}{Full-Stack \& DevOps Developer Intern}{Java, Spring Boot, Angular, AWS}
\begin{itemize}
  \item Developed new features for an internal rights, contracts and film catalogue platform used by \textbf{100+ employees}.
  \item Set up the GitLab CI/CD pipeline for a new module; deployed across three environments (DEV, UAT, PROD).
  \item Worked on the AWS infrastructure (EC2, RDS, IAM) within an Agile team.
\end{itemize}

\vspace{3pt}
\entry{Hitachi Solutions France}{June -- July 2022}{IT Consultant Intern}{Python, Django}
\begin{itemize}
  \item Prototyped a supplier, product and contract management application for Sonepar, following Agile SAFe.
\end{itemize}

%% ══ SKILLS ═══════════════════════════════════════════════════════════════════
\section{Skills}

\begin{tabularx}{\textwidth}{@{}l@{\hspace{1.2em}}X@{}}
  \textbf{Computer Vision \& AI} & Image processing, deep learning, PyTorch, OpenCV, YOLO, Florence-2, XGBoost \\
  \textbf{GPU \& Graphics}       & CUDA, Nsight, OpenGL, GLSL, compute shaders \\
  \textbf{Programming}           & C, C++, Python, Java, TypeScript, SQL \\
  \textbf{Frameworks}            & Spring Boot, Quarkus, FastAPI, Django, Angular, React \\
  \textbf{Tools}                 & Linux, Git, Docker, GitLab CI, AWS, PostgreSQL, MongoDB, SQLite \\
  \textbf{Languages}             & French (native), English (professional, TOEIC 845)
\end{tabularx}

\end{document}
